% =====================================================================
%  Цель и задание лабораторной работы.
% =====================================================================

\textbf{Цель:} изучить механизм списков контроля доступа (ACL) и получить
практические навыки настройки прав доступа к объектам файловой системы
Linux.

\vspace{1em}

\textbf{Задание}:
\begin{enumerate}
  \item Создать группы пользователей \texttt{group\_iit1}, \texttt{group\_iit2};
  создать пользователей \texttt{iit11}, \texttt{iit12} (группа \texttt{group\_iit1})
  и \texttt{iit21}, \texttt{iit22} (группа \texttt{group\_iit2}); предоставить
  \texttt{iit21} административные привилегии; создать пользователя \texttt{iit3}.
  \item Создать каталог \texttt{pzs} и вложенные каталоги \texttt{pzs11}--\texttt{pzs15}
  с правами чтения, записи и выполнения только для владельца, только для группы,
  только для остальных пользователей, для всех пользователей и только для
  администратора соответственно.
  \item От имени пользователя \texttt{iit11} создать в каталогах
  \texttt{pzs11}--\texttt{pzs14} файлы, реализующие все сочетания прав (чтение,
  запись, выполнение и их комбинации) для каждого из четырёх классов субъектов
  доступа.
  \item Для каждого созданного файла проверить возможность чтения, редактирования
  и запуска пользователями \texttt{iit11}, \texttt{iit12}, \texttt{iit21},
  \texttt{iit22}, \texttt{iit3} и суперпользователем.
  \item Запустить от имени \texttt{iit11} файлы, предназначенные только для
  выполнения, и проверить, кто из пользователей способен остановить запущенный
  процесс.
  \item Для каждого каталога проверить возможность просмотра содержимого,
  создания и удаления файлов всеми пользователями.
  \item Удалить созданные файлы, каталоги, пользователей и группы.
\end{enumerate}
